% Documentary recovery for Article V; RESULTADO_RECUPERADO.
% Source owner: /Users/ruben/Documents/New project/output/REVISION_EDITORIAL_HMT_MD_20260905/01_FUENTES/INTEGRAL/tree/New project/output/ENSAMBLADO_CAUSAL_RESTAURADO_HMT_MD_20260905/fuente/manuscrito/integracion_83/parte_iv/body/B011_ch53_trayectorias_cuanticas_6b08e0e532c9_04ab_postulados_cuanticos_genealogicos_rev2.tex
% Source lines: 4--95, 178--206.
% Only reference labels and, where appropriate, heading levels are adapted.
% The recovery matrix preserves the correspondence and dependencies.
\section{Reconstruction of quantum principles from the genealogical state}
\label{v:rec:sec:postulados-cuanticos-genealogicos-rev2}
\index{quantum mechanics!genealogical reconstruction}
\index{state!Hilbert realization}

The Hilbert-space formulation organizes quantum mechanics in terms of states,
self-adjoint observables, unitary evolution, tensor composition, and a
probability rule for outcomes~\cite{vonNeumann1932}. HMT--MD
constructs a layer prior to that representation: a compatible section,
its observable route, the memory of transport, and the reader that expresses an
outcome. Hilbert space is not replaced; it appears as the codomain of
an explicitly specified linear realization.

The typed sequence is
\[
 \boxed{\begin{aligned}
 \text{genealogical section}
 &\longrightarrow \text{linear realization}
 \longrightarrow \text{phase and superposition}\\
 &\longrightarrow \text{observable and evolution}
 \longrightarrow \text{probability and measurement}
 \end{aligned}}
\]
Each arrow preserves its domain. In particular, the state vector does not
replace the TPK section, just as the real value does not replace
the numerical genealogy that produces it.

\subsection{Genealogical state and Hilbert realization}
\label{v:rec:realizacion_evolucion:section:02}

\begin{definition}[Realized genealogical state]
\label{v:rec:realizacion_evolucion:statement:01}
A genealogical state is a quintuple
\begin{equation}
 \mathfrak G=(x_\infty,\Gamma,\mathcal F,\mathcal M,\mathcal R),
 \label{v:rec:eq:estado-genealogico-realizado-rev2}
\end{equation}
where \(x_\infty\in\Omega_\infty\) is a compatible section,
\(\Gamma\) is its observable route, \(\mathcal F\) is the collection of realizations,
\(\mathcal M\) is the memory of the lifts, and \(\mathcal R\) is the specified
reader. Let \(\mathsf P_{\rm TPK}\) be the category whose objects are
prepared TPK sections, finite or coinductive, and whose morphisms are
recorded routes compatible with their boundary data. In particular,
depth restriction is a morphism
\[
 r_{n+1,n}:x_{n+1}\longrightarrow x_n.
\]
Let \(\mathsf{Hilb}_{\mathbb R,J}\) be the category whose objects
are pairs \((\mathcal H_{\mathbb R},J)\), with \(\mathcal H_{\mathbb R}\)
a real Hilbert space and \(J\) orthogonal, \(J^2=-I\), and whose morphisms are
bounded real-linear maps \(T\) intertwining the complex
structures, \(TJ_1=J_2T\). A quantum realization is a functor
\[
 \mathcal Q:\mathsf P_{\rm TPK}\longrightarrow
 \mathsf{Hilb}_{\mathbb R,J},
\]
and a \emph{pointed realization} is the pair \((\mathcal Q,\Psi)\) assigning
to each prepared section \(x\) a nonzero vector
\(\Psi_x\in\mathcal Q(x)\), compatible with the restrictions of the inverse
system.
\end{definition}

With the convention of linearity in the second argument, the associated Hermitian form
is written
\[
 \langle u,v\rangle_{\mathbb C}
 =\langle u,v\rangle_{\mathbb R}
  -i\langle u,J_E v\rangle_{\mathbb R}.
\]
Choosing the chart \(J_E\leftrightarrow i\) realizes the complex structure that
the formal kernel and its development construct as an internal square root of \(-I\). For prefixes of
depth \(n\), set
\(\mathcal H_n:=\mathcal Q(x_n)\) and
\[
 p_{n+1,n}:=\mathcal Q(r_{n+1,n}):
 \mathcal H_{n+1}\longrightarrow\mathcal H_n.
\]
Compatibility between restriction and realization requires
\begin{equation}
 p_{n+1,n}\Psi_{n+1}(x_{n+1})=\Psi_n(x_n).
 \label{v:rec:eq:compatibilidad-realizacion-hilbert-rev2}
\end{equation}
If \(r_{\infty,n}:x_\infty\to x_n\) is the coinductive restriction and
\(p_n:=\mathcal Q(r_{\infty,n}):\mathcal H_\infty\to\mathcal H_n\), a
limit \(\mathcal H_\infty\) realizes the complete section when
\(p_n\Psi_\infty(x_\infty)=\Psi_n(x_n)\) at every depth.

\begin{principle}[Realized quantum state]
\label{v:rec:realizacion_evolucion:statement:02}
The pure state associated with a section is represented by the ray of the
nonzero vector \(\Psi=\Psi_\infty(x_\infty)\); a statistical preparation is
represented by a positive, trace-one density operator on the same
realization. Global-phase identification acts in the representation:
the section and its record remain in the genealogical domain.
\end{principle}

\subsection{Evolution, Hamiltonian, and Schrödinger equation}
\label{v:rec:realizacion_evolucion:section:03}

Genealogical evolution preserves composition of transports. Its
realization satisfies
\[
 U(t_3,t_2)U(t_2,t_1)=U(t_3,t_1),
 \qquad U(t,t)=I.
\]
When it preserves the inner product and satisfies
\(U(t)J_E=J_EU(t)\), it is orthogonal in the real chart and unitary in the
complex chart. For a strongly continuous one-parameter unitary group,
Stone's theorem provides a self-adjoint, \(J_E\)-linear generator
\(H\)~\cite{Stone1932}:
\[
 U(t)=\exp\!\left(-J_E H t/\hbar_{\rm int}\right).
\]
On the domain \(D(H)\), differentiation with respect to time gives
\begin{equation}
 \hbar_{\rm int}J_E\dot\Psi(t)=H\Psi(t),
 \label{v:rec:eq:schrodinger-genealogico-rev2}
\end{equation}
and the chart \(J_E\leftrightarrow i\) yields
\(i\hbar_{\rm int}\dot\Psi=H\Psi\).

The calendar \(C_{108}\) realizes this correspondence in finite form:
the self-adjoint Hamiltonian constructed for it has, over \(t_0\),
the unitary one-iteration shift as its exponential. The
continuous group additionally requires strong continuity; it does not follow
from the periodicity of the calendar alone.
